\section{报告定位与交付边界}

\begin{reportbox}{核心定位}
本报告是 SciScope 项目面向赛题交付的独立数据分析报告, 目标是回答“公开科技文献数据本身揭示了什么”。它与系统设计 PDF 并列: 设计 PDF 解释智能体架构和工程路线, 本报告沉淀文献分布、关键词演化、作者合作网络、趋势预测和可复现分析结论。
\end{reportbox}

\subsection{当前分析阶段}

当前版本已经从报告底座推进到 5 万级数据分析版本。分析资产覆盖 50,498 条 raw records、49,067 条去重语料记录、40,137 条近五年智能体语料记录, 并生成文献分布、关键词演化、作者合作网络和数据质量审计图表。

\begin{center}
\begin{tabularx}{\textwidth}{>{\bfseries}p{34mm}X X}
\toprule
交付维度 & 当前已有内容 & 下一轮增强内容 \\
\midrule
数据来源 & OpenAlex、arXiv、PubMed、PMC、Crossref、DOAJ 六源采集 & Semantic Scholar、CORE 等 API-key 增强源 \\
数据规模 & 50,498 条 raw records 与 49,067 条去重语料记录 & DOI/title/year 深度去重与字段补全 \\
分析主题 & 分布、关键词、作者网络、质量审计、智能体映射 & 动态主题模型、趋势预测、图谱社区发现 \\
报告产物 & 多章节 LaTeX、可编译 PDF、13 张图表资产 & 评审版摘要页与前端联动展示 \\
\bottomrule
\end{tabularx}
\end{center}

\subsection{报告读者}

\begin{itemize}
  \item \textbf{评委}: 需要快速看见数据真实、分析深入、结论可解释。
  \item \textbf{科研用户}: 需要理解某一领域文献热点、演化路径和关键合作网络。
  \item \textbf{系统开发者}: 需要把报告中的统计资产复用到 RAG、推荐、趋势和知识图谱模块。
\end{itemize}

\clearpage
